\documentclass[10pt]{article}
\PassOptionsToPackage{obeyspaces,spaces}{url} % Los comandos con \path conservan sus espacios
% Estilo común del SRS, el SAD y el plan de pruebas. Reproduce el formato del SRS y el
% SAD originales: carta, márgenes de una pulgada, Times de 10 puntos en el cuerpo.
% Las rutas son relativas a la carpeta del documento: se compila desde ella
% (latexmk -cd) y se carga con % Estilo común del SRS, el SAD y el plan de pruebas. Reproduce el formato del SRS y el
% SAD originales: carta, márgenes de una pulgada, Times de 10 puntos en el cuerpo.
% Las rutas son relativas a la carpeta del documento: se compila desde ella
% (latexmk -cd) y se carga con % Estilo común del SRS, el SAD y el plan de pruebas. Reproduce el formato del SRS y el
% SAD originales: carta, márgenes de una pulgada, Times de 10 puntos en el cuerpo.
% Las rutas son relativas a la carpeta del documento: se compila desde ella
% (latexmk -cd) y se carga con \input{../comun/estilo}.

\usepackage[T1]{fontenc}
\usepackage[utf8]{inputenc}
\usepackage[spanish,es-tabla,es-noshorthands,es-lcroman]{babel}
\usepackage{tgtermes}
\usepackage{tgheros}
\usepackage[letterpaper,margin=1in]{geometry}
\usepackage{graphicx}
\usepackage{array}
\usepackage{tabularx}
\usepackage{setspace}
\usepackage{titlesec}
\usepackage{titletoc}
\usepackage{caption}
\usepackage[hidelinks]{hyperref}

\setstretch{1.15}
\setlength{\parindent}{0pt}
\setlength{\parskip}{6pt plus 2pt}

\addto\captionsspanish{%
  \renewcommand{\contentsname}{Tabla de Contenidos}%
  \renewcommand{\listtablename}{Lista de Tablas}%
  \renewcommand{\listfigurename}{Lista de Figuras}%
}

\titleformat{\section}{\normalfont\bfseries\fontsize{12}{14}\selectfont}{\thesection.}{0.3em}{}
\titleformat{\subsection}{\normalfont\bfseries\fontsize{11}{13}\selectfont}{\thesubsection}{0.3em}{}
\titleformat{\subsubsection}{\normalfont\bfseries\normalsize}{\thesubsubsection}{0.3em}{}
\titlespacing*{\section}{0pt}{14pt}{8pt}
\titlespacing*{\subsection}{0pt}{12pt}{6pt}
\titlespacing*{\subsubsection}{0pt}{10pt}{4pt}

\titlecontents{section}[0pt]{\addvspace{2pt}\bfseries}
  {\thecontentslabel.\ }{}{\titlerule*[0.3em]{.}\contentspage}
\titlecontents{subsection}[1.5em]{}
  {\thecontentslabel\ }{}{\titlerule*[0.3em]{.}\contentspage}
\titlecontents{subsubsection}[3em]{}
  {\thecontentslabel\ }{}{\titlerule*[0.3em]{.}\contentspage}
\titlecontents{table}[0.75in]{}
  {\tablename~\thecontentslabel.\ }{}{\titlerule*[0.3em]{.}\contentspage}
\titlecontents{figure}[0.75in]{}
  {\figurename~\thecontentslabel.\ }{}{\titlerule*[0.3em]{.}\contentspage}

\captionsetup{labelsep=period,font=it,justification=centering}

\newcommand{\kpNombre}{}
\newcommand{\kpVersion}{}
\newcommand{\kpFecha}{}
\newcommand{\NombreDocumento}[1]{\renewcommand{\kpNombre}{#1}}
\newcommand{\VersionDocumento}[1]{\renewcommand{\kpVersion}{#1}}
\newcommand{\FechaDocumento}[1]{\renewcommand{\kpFecha}{#1}}

\newcommand{\Portada}{%
  \hypersetup{
    pdftitle={\kpNombre\ - Kit Pagos Colombia},
    pdfauthor={Andres Henao Niño, Joan Emmanuel Orduz Chia, Joshua Prieto Zambrano, David Estevan Rodriguez Jurado}
  }%
  \begin{titlepage}
    \centering
    \setstretch{1}%
    \setlength{\parskip}{0pt}%
    \fontsize{12}{16}\selectfont
    \includegraphics[height=6.96cm]{../comun/logo-javeriana}\par
    \vspace{46pt}
    {\fontsize{14}{16}\selectfont\bfseries\kpNombre\par}
    \vspace{16pt}
    Presentan\par
    \vspace{16pt}
    Andres Henao Niño\par
    Joan Emmanuel Orduz Chia\par
    Joshua Prieto Zambrano\par
    David Estevan Rodriguez Jurado\par
    \vspace{31pt}
    Director\par
    Jaime Andrés Pavlich Mariscal\par
    \vspace{21pt}
    Pontificia Universidad Javeriana\par
    Facultad de Ingeniería\par
    Trabajo de Grado\par
    \vspace{16pt}
    Bogotá D.C.\par
    \kpFecha\par
    Versión \kpVersion\par
  \end{titlepage}%
  \pagenumbering{roman}%
}

% Cada fila se agrega con \Cambio{versión}{fecha}{descripción}{autor(es)}.
\newenvironment{HistorialDeCambios}{%
  \section*{Historial de Cambios}%
  \renewcommand{\arraystretch}{1.3}%
  \tabularx{\textwidth}{|c|c|X|>{\raggedright\arraybackslash}p{4.2cm}|}
  \hline
  \textbf{Versión} & \textbf{Fecha} & \textbf{Descripción del cambio} & \textbf{Autor(es)} \\
  \hline
}{%
  \endtabularx
  \clearpage
}
\newcommand{\Cambio}[4]{#1 & #2 & #3 & #4 \\ \hline}

% La plantilla PTP pide que la numeración arábiga empiece en la Introducción y que las
% listas de tablas y figuras no entren en la tabla de contenidos.
\newcommand{\Indices}{%
  \tableofcontents
  \clearpage
  \listoftables
  \listoffigures
  \clearpage
  \pagenumbering{arabic}%
}.

\usepackage[T1]{fontenc}
\usepackage[utf8]{inputenc}
\usepackage[spanish,es-tabla,es-noshorthands,es-lcroman]{babel}
\usepackage{tgtermes}
\usepackage{tgheros}
\usepackage[letterpaper,margin=1in]{geometry}
\usepackage{graphicx}
\usepackage{array}
\usepackage{tabularx}
\usepackage{setspace}
\usepackage{titlesec}
\usepackage{titletoc}
\usepackage{caption}
\usepackage[hidelinks]{hyperref}

\setstretch{1.15}
\setlength{\parindent}{0pt}
\setlength{\parskip}{6pt plus 2pt}

\addto\captionsspanish{%
  \renewcommand{\contentsname}{Tabla de Contenidos}%
  \renewcommand{\listtablename}{Lista de Tablas}%
  \renewcommand{\listfigurename}{Lista de Figuras}%
}

\titleformat{\section}{\normalfont\bfseries\fontsize{12}{14}\selectfont}{\thesection.}{0.3em}{}
\titleformat{\subsection}{\normalfont\bfseries\fontsize{11}{13}\selectfont}{\thesubsection}{0.3em}{}
\titleformat{\subsubsection}{\normalfont\bfseries\normalsize}{\thesubsubsection}{0.3em}{}
\titlespacing*{\section}{0pt}{14pt}{8pt}
\titlespacing*{\subsection}{0pt}{12pt}{6pt}
\titlespacing*{\subsubsection}{0pt}{10pt}{4pt}

\titlecontents{section}[0pt]{\addvspace{2pt}\bfseries}
  {\thecontentslabel.\ }{}{\titlerule*[0.3em]{.}\contentspage}
\titlecontents{subsection}[1.5em]{}
  {\thecontentslabel\ }{}{\titlerule*[0.3em]{.}\contentspage}
\titlecontents{subsubsection}[3em]{}
  {\thecontentslabel\ }{}{\titlerule*[0.3em]{.}\contentspage}
\titlecontents{table}[0.75in]{}
  {\tablename~\thecontentslabel.\ }{}{\titlerule*[0.3em]{.}\contentspage}
\titlecontents{figure}[0.75in]{}
  {\figurename~\thecontentslabel.\ }{}{\titlerule*[0.3em]{.}\contentspage}

\captionsetup{labelsep=period,font=it,justification=centering}

\newcommand{\kpNombre}{}
\newcommand{\kpVersion}{}
\newcommand{\kpFecha}{}
\newcommand{\NombreDocumento}[1]{\renewcommand{\kpNombre}{#1}}
\newcommand{\VersionDocumento}[1]{\renewcommand{\kpVersion}{#1}}
\newcommand{\FechaDocumento}[1]{\renewcommand{\kpFecha}{#1}}

\newcommand{\Portada}{%
  \hypersetup{
    pdftitle={\kpNombre\ - Kit Pagos Colombia},
    pdfauthor={Andres Henao Niño, Joan Emmanuel Orduz Chia, Joshua Prieto Zambrano, David Estevan Rodriguez Jurado}
  }%
  \begin{titlepage}
    \centering
    \setstretch{1}%
    \setlength{\parskip}{0pt}%
    \fontsize{12}{16}\selectfont
    \includegraphics[height=6.96cm]{../comun/logo-javeriana}\par
    \vspace{46pt}
    {\fontsize{14}{16}\selectfont\bfseries\kpNombre\par}
    \vspace{16pt}
    Presentan\par
    \vspace{16pt}
    Andres Henao Niño\par
    Joan Emmanuel Orduz Chia\par
    Joshua Prieto Zambrano\par
    David Estevan Rodriguez Jurado\par
    \vspace{31pt}
    Director\par
    Jaime Andrés Pavlich Mariscal\par
    \vspace{21pt}
    Pontificia Universidad Javeriana\par
    Facultad de Ingeniería\par
    Trabajo de Grado\par
    \vspace{16pt}
    Bogotá D.C.\par
    \kpFecha\par
    Versión \kpVersion\par
  \end{titlepage}%
  \pagenumbering{roman}%
}

% Cada fila se agrega con \Cambio{versión}{fecha}{descripción}{autor(es)}.
\newenvironment{HistorialDeCambios}{%
  \section*{Historial de Cambios}%
  \renewcommand{\arraystretch}{1.3}%
  \tabularx{\textwidth}{|c|c|X|>{\raggedright\arraybackslash}p{4.2cm}|}
  \hline
  \textbf{Versión} & \textbf{Fecha} & \textbf{Descripción del cambio} & \textbf{Autor(es)} \\
  \hline
}{%
  \endtabularx
  \clearpage
}
\newcommand{\Cambio}[4]{#1 & #2 & #3 & #4 \\ \hline}

% La plantilla PTP pide que la numeración arábiga empiece en la Introducción y que las
% listas de tablas y figuras no entren en la tabla de contenidos.
\newcommand{\Indices}{%
  \tableofcontents
  \clearpage
  \listoftables
  \listoffigures
  \clearpage
  \pagenumbering{arabic}%
}.

\usepackage[T1]{fontenc}
\usepackage[utf8]{inputenc}
\usepackage[spanish,es-tabla,es-noshorthands,es-lcroman]{babel}
\usepackage{tgtermes}
\usepackage{tgheros}
\usepackage[letterpaper,margin=1in]{geometry}
\usepackage{graphicx}
\usepackage{array}
\usepackage{tabularx}
\usepackage{setspace}
\usepackage{titlesec}
\usepackage{titletoc}
\usepackage{caption}
\usepackage[hidelinks]{hyperref}

\setstretch{1.15}
\setlength{\parindent}{0pt}
\setlength{\parskip}{6pt plus 2pt}

\addto\captionsspanish{%
  \renewcommand{\contentsname}{Tabla de Contenidos}%
  \renewcommand{\listtablename}{Lista de Tablas}%
  \renewcommand{\listfigurename}{Lista de Figuras}%
}

\titleformat{\section}{\normalfont\bfseries\fontsize{12}{14}\selectfont}{\thesection.}{0.3em}{}
\titleformat{\subsection}{\normalfont\bfseries\fontsize{11}{13}\selectfont}{\thesubsection}{0.3em}{}
\titleformat{\subsubsection}{\normalfont\bfseries\normalsize}{\thesubsubsection}{0.3em}{}
\titlespacing*{\section}{0pt}{14pt}{8pt}
\titlespacing*{\subsection}{0pt}{12pt}{6pt}
\titlespacing*{\subsubsection}{0pt}{10pt}{4pt}

\titlecontents{section}[0pt]{\addvspace{2pt}\bfseries}
  {\thecontentslabel.\ }{}{\titlerule*[0.3em]{.}\contentspage}
\titlecontents{subsection}[1.5em]{}
  {\thecontentslabel\ }{}{\titlerule*[0.3em]{.}\contentspage}
\titlecontents{subsubsection}[3em]{}
  {\thecontentslabel\ }{}{\titlerule*[0.3em]{.}\contentspage}
\titlecontents{table}[0.75in]{}
  {\tablename~\thecontentslabel.\ }{}{\titlerule*[0.3em]{.}\contentspage}
\titlecontents{figure}[0.75in]{}
  {\figurename~\thecontentslabel.\ }{}{\titlerule*[0.3em]{.}\contentspage}

\captionsetup{labelsep=period,font=it,justification=centering}

\newcommand{\kpNombre}{}
\newcommand{\kpVersion}{}
\newcommand{\kpFecha}{}
\newcommand{\NombreDocumento}[1]{\renewcommand{\kpNombre}{#1}}
\newcommand{\VersionDocumento}[1]{\renewcommand{\kpVersion}{#1}}
\newcommand{\FechaDocumento}[1]{\renewcommand{\kpFecha}{#1}}

\newcommand{\Portada}{%
  \hypersetup{
    pdftitle={\kpNombre\ - Kit Pagos Colombia},
    pdfauthor={Andres Henao Niño, Joan Emmanuel Orduz Chia, Joshua Prieto Zambrano, David Estevan Rodriguez Jurado}
  }%
  \begin{titlepage}
    \centering
    \setstretch{1}%
    \setlength{\parskip}{0pt}%
    \fontsize{12}{16}\selectfont
    \includegraphics[height=6.96cm]{../comun/logo-javeriana}\par
    \vspace{46pt}
    {\fontsize{14}{16}\selectfont\bfseries\kpNombre\par}
    \vspace{16pt}
    Presentan\par
    \vspace{16pt}
    Andres Henao Niño\par
    Joan Emmanuel Orduz Chia\par
    Joshua Prieto Zambrano\par
    David Estevan Rodriguez Jurado\par
    \vspace{31pt}
    Director\par
    Jaime Andrés Pavlich Mariscal\par
    \vspace{21pt}
    Pontificia Universidad Javeriana\par
    Facultad de Ingeniería\par
    Trabajo de Grado\par
    \vspace{16pt}
    Bogotá D.C.\par
    \kpFecha\par
    Versión \kpVersion\par
  \end{titlepage}%
  \pagenumbering{roman}%
}

% Cada fila se agrega con \Cambio{versión}{fecha}{descripción}{autor(es)}.
\newenvironment{HistorialDeCambios}{%
  \section*{Historial de Cambios}%
  \renewcommand{\arraystretch}{1.3}%
  \tabularx{\textwidth}{|c|c|X|>{\raggedright\arraybackslash}p{4.2cm}|}
  \hline
  \textbf{Versión} & \textbf{Fecha} & \textbf{Descripción del cambio} & \textbf{Autor(es)} \\
  \hline
}{%
  \endtabularx
  \clearpage
}
\newcommand{\Cambio}[4]{#1 & #2 & #3 & #4 \\ \hline}

% La plantilla PTP pide que la numeración arábiga empiece en la Introducción y que las
% listas de tablas y figuras no entren en la tabla de contenidos.
\newcommand{\Indices}{%
  \tableofcontents
  \clearpage
  \listoftables
  \listoffigures
  \clearpage
  \pagenumbering{arabic}%
}
\usepackage{float} % Permite usar [H] para posicionamiento exacto de tablas
\usepackage{longtable} % Tablas que pueden partirse entre páginas
\usepackage{xltabular} % Tablas tabularx que pueden partirse entre páginas
\usepackage{xurl} % Saltos de línea en rutas, URL e identificadores largos

\NombreDocumento{Plan de Pruebas del Proyecto (PTP)}
\VersionDocumento{0.4}
\FechaDocumento{noviembre de 2026}

\begin{document}

\Portada

\begin{HistorialDeCambios}
  \Cambio{0.1}{28/10/2026}{Portada, historial de cambios e índices con el estilo común.}{Joan Emmanuel Orduz Chia}
  \Cambio{0.2}{09/11/2026}{Ítems de prueba, estrategia, siete casos diseñados y estructura de resultados.}{Andres Henao Niño, Joan Emmanuel Orduz Chia, Joshua Prieto Zambrano, David Estevan Rodriguez Jurado}
  \Cambio{0.3}{09/11/2026}{Estrategia de carga aislada de las pasarelas, desglose por archivo de las pruebas unitarias y de integración, CP-02 a CP-07 en los anexos y correcciones de la revisión.}{Joan Emmanuel Orduz Chia}
  \Cambio{0.4}{11/11/2026}{Resultados de la ejecución: los siete casos, las suites, las pruebas de contrato, la carga, los defectos con su corrección y las evidencias.}{Joan Emmanuel Orduz Chia}
\end{HistorialDeCambios}

\Indices

\section{Introducción}

Este documento es el Plan de Pruebas del Proyecto (PTP) de Kit Pagos Colombia. Describe qué se va a probar, con qué estrategia, en qué ambientes y con qué criterios de terminación, y diseña los casos que se ejecutarán para verificar el comportamiento del framework. No sustituye al SRS ni al documento de arquitectura: los toma como fuente y convierte sus requisitos en verificaciones concretas.

\subsection{Propósito}

El propósito del documento es que cualquier integrante del equipo, o un evaluador externo, pueda responder cuatro preguntas sin consultar otro archivo:

\begin{itemize}
  \item \textbf{Qué se prueba.} Los componentes del framework y los productos que lo consumen, delimitados en la Sección~\ref{sec:items}.
  \item \textbf{Cómo se prueba.} Los tipos de prueba, las herramientas, los ambientes y los criterios de completitud de cada tipo, descritos en la Sección~\ref{sec:estrategia}.
  \item \textbf{Con qué casos.} Siete casos diseñados, uno funcional completo en el cuerpo y los seis restantes en los anexos, con sus pasos, datos y resultados esperados (Sección~\ref{sec:casos} y Anexos).
  \item \textbf{Cómo se reporta.} La estructura de resultados, el registro de defectos y las métricas que se llenarán al ejecutar las pruebas (Sección~\ref{sec:resultados}).
\end{itemize}

\subsection{Alcance}

El alcance cubre los tres componentes del artefacto y los cuatro prototipos de la Fase~5:

\begin{itemize}
  \item El \textbf{SDK de servidor} \texttt{kit-pagos-colombia}, que unifica las cuatro pasarelas colombianas.
  \item El módulo de navegador \texttt{kit-pagos-colombia/browser}, que tokeniza la tarjeta sin que el número pase por el servidor del comercio.
  \item La \textbf{API de Simulación}, tanto las rutas que imitan a las pasarelas bajo \texttt{/v1/sim} como la capa REST propia bajo \texttt{/v1/api}.
  \item Los prototipos A (\texttt{checkout-con-sdk-solo-back}), B (\texttt{checkout-con-sdk-front-back}), C (\texttt{checkout-frontend-rest}) y D (\texttt{checkout-sin-solucion}).
\end{itemize}

Quedan fuera del alcance las operaciones CRUD de la tienda, la seguridad de la aplicación anfitriona, la ayuda en línea y el rendimiento bajo carga. Tampoco se prueba el cumplimiento PCI~DSS de la página del comercio, que depende de la aplicación que integra el framework. Las pruebas de contrato contra los sandboxes reales se diseñan y se cuentan, pero su ejecución requiere credenciales y red, y por eso se reporta por separado.

\subsection{Audiencia y convenciones}

El documento está dirigido al equipo de desarrollo, al director del trabajo de grado y a los evaluadores. Los comandos se escriben desde la raíz del repositorio salvo que se indique otra carpeta. Los nombres de archivos, rutas, clases y variables se conservan en su forma original en inglés o en el idioma del código; las explicaciones están en español.

Cada caso de prueba declara de forma explícita su \textbf{ambiente}, porque el resultado de una prueba de pagos depende de dónde se ejecute: no es lo mismo el simulador local, que es determinista, que un sandbox real, que puede declinar cuando quiere. Esa distinción se mantiene en todo el documento y es la razón por la que los casos funcionales se ejecutan contra el simulador y las pruebas de contrato se reportan aparte.

\subsection{Estructura del documento}

\begin{description}
  \item[Sección~\ref{sec:items}.] Ítems de prueba: los componentes, su descripción y el objetivo de probar cada uno.
  \item[Sección~\ref{sec:estrategia}.] Estrategia: tipos de prueba, herramientas, ambientes, datos y criterios de completitud.
  \item[Sección~\ref{sec:casos}.] Casos de prueba: el resumen de los siete casos y el caso CP-01 completo.
  \item[Sección~\ref{sec:resultados}.] Resultados: la estructura que se llenará en el issue \#174, con la escala de estados, el registro de defectos y las métricas.
  \item[Sección~\ref{sec:referencias}.] Referencias bibliográficas en formato IEEE.
  \item[Anexos.] Los casos CP-02 a CP-07 completos, la matriz de trazabilidad, los comandos de la corrida y el glosario.
\end{description}

\subsection{Relación con la ejecución}

Este documento diseña las pruebas y deja preparada la estructura de resultados. La ejecución, las evidencias y el registro de defectos corresponden al issue \#174. Por esa razón, las columnas de resultado obtenido, estado y notas de las tablas de pasos aparecen vacías: se llenarán con la corrida real y no con una expectativa.

\section{Ítems de prueba}
\textbf{Propósito:} Que el equipo de desarrollo sepa cuáles son los elementos del proyecto de software que se van a probar.

\textbf{Contenido:}
\begin{itemize}
  \item Tabla con los ítems que se van a probar en el ámbito de este documento. Tipos de ítem pueden incluir: métodos, clases, componentes, servicios WEB y pantallas de formularios, entre otros.
  \item Incluir nombre, descripción y cuál será el objetivo de probar ese ítem.
\end{itemize}

\textbf{Tamaño sugerido:} 1 página.

\begin{table}[H]
  \centering
  \small
  \renewcommand{\arraystretch}{1.3}
  \setlength{\tabcolsep}{4pt}
  \begin{tabularx}{\textwidth}{|>{\raggedright\arraybackslash}p{3cm}|>{\raggedright\arraybackslash}X|>{\raggedright\arraybackslash}X|}
    \hline
    \textbf{Nombre del Ítem} & \textbf{Descripción} & \textbf{Objetivo de la prueba} \\ \hline
    [Nombre] & [Descripción del elemento] & [Objetivo de probar el ítem] \\ \hline
  \end{tabularx}
  \caption{Tabla de ítems de prueba}
  \label{tab:items-prueba}
\end{table}
\section{Estrategia de pruebas}
\textbf{Propósito:} Que el equipo de desarrollo sepa cuáles son los tipos de pruebas que se van a realizar y la estrategia general para llevar a cabo esas pruebas.

\textbf{Contenido:}
\begin{itemize}
  \item Explicar los tipos de pruebas (Funcionales y no funcionales) que se van a realizar a los distintos elementos que estarán bajo pruebas. Ejemplo: Pruebas unitarias, de integración, de sistema, de regresión, de rendimiento y de aceptación del usuario. Debe incluir al menos dos tipos de pruebas en su Plan de Pruebas.
  \item Plantee la estrategia para realizar cada tipo de prueba. Especifique la(s) herramienta, ambiente, criterios de completitud y otra información que considere relevante. Incluya referencias a las herramientas utilizadas (e.g. JUnit, Selenium, Mockito, Bugzilla, etc.).
\end{itemize}

\textbf{Tamaño sugerido:} 3 páginas.
\section{Casos de prueba}
\textbf{Propósito:} Desarrollar los casos de pruebas funcionales y no funcionales para el proyecto.

\textbf{Contenido:}
\begin{itemize}
  \item Diseñe al menos 5 casos de pruebas para funcionalidades clave de su software.
  \item No incluir CRUDs, seguridad o ayuda.
  \item Muestre un ejemplo de un caso de prueba. Los demás pueden estar en un anexo.
  \item Diseñe al menos una prueba para un atributo de calidad.
  \item Puede utilizar la plantilla de la Figura~\ref{fig:plantilla-casos-prueba} para los casos de prueba u otra similar.
\end{itemize}

\begin{table}[H]
  \centering
  \small
  \renewcommand{\arraystretch}{1.2}
  \setlength{\tabcolsep}{3pt}
  \begin{tabularx}{\textwidth}{|>{\raggedright\arraybackslash}X|>{\raggedright\arraybackslash}X|}
    \hline
    \textbf{Id:} & \textbf{Author:} \\ \hline
    \textbf{Date created:} & \textbf{Title:} \\ \hline
    \textbf{Tester:} & \textbf{Date tested:} \\ \hline
    \multicolumn{2}{|l|}{\textbf{Description:}} \\ \hline
    \multicolumn{2}{|l|}{\textbf{Type:} (Functional, Usability, Performance, Reliability, Maintainability, Security, etc.)} \\ \hline
    \multicolumn{2}{|l|}{\textbf{Preconditions:}} \\ \hline
  \end{tabularx}
  \vspace{3pt}

  \setlength{\tabcolsep}{2pt}
  \begin{tabularx}{\textwidth}{|c|>{\raggedright\arraybackslash}X|>{\raggedright\arraybackslash}p{1.1cm}|>{\raggedright\arraybackslash}X|>{\raggedright\arraybackslash}X|>{\raggedright\arraybackslash}p{1.4cm}|>{\raggedright\arraybackslash}p{1.1cm}|}
    \hline
    \textbf{Test} & \textbf{Steps} & \textbf{Data} & \textbf{Expected Result} & \textbf{Actual Result} & \textbf{Status} & \textbf{Notes} \\ \hline
    1 & Invoke add operation & 5, 2 & 7 & 7 & Passed/Failed & \\ \hline
    2 & Invoke operation with negative values & -9, 3 & -6 & & & \\ \hline
    3 & Invoke 10000 times the add operation using a load testing tool & 5, 2 & All operations are processed & & & \\ \hline
    4 & \dots & & & & & \\ \hline
  \end{tabularx}
  \caption{Plantilla de casos de prueba}
  \label{fig:plantilla-casos-prueba}
\end{table}
\section{Resultados}
\textbf{Propósito:} Mostrar las evidencias de la aplicación del proceso de pruebas en el proyecto de software.

\textbf{Contenido:}
\begin{itemize}
  \item Para cada prueba diseñada en este documento, mostrar las evidencias de la ejecución de cada una de las pruebas planificadas (puede usar la tabla de resultados esperados de la Figura 1 o un formato similar).
  \item Para los casos en que la prueba falle, debe incluir un comentario que explique cuál es el error encontrado.
\end{itemize}
\section{Referencias}
\label{sec:referencias}

\begin{enumerate}
  \renewcommand{\labelenumi}{[\arabic{enumi}]}

  \item Kit Pagos Colombia, ``Repositorio del proyecto,'' GitHub. [En línea]. Disponible en: \url{https://github.com/PurosBrothers/Kit-Pagos-Colombia---Tesis}

  \item Kit Pagos Colombia, ``Especificación de Requisitos de Software (SRS), versión 2.0,'' \path{docs/final-doc/srs/}, 2026.

  \item Kit Pagos Colombia, ``Documento de Arquitectura de Software (SAD),'' \path{docs/final-doc/sad/}, 2026.

  \item Kit Pagos Colombia, ``API de Simulación: el segundo componente,'' \path{docs/02-arquitectura/3-api-de-simulacion.md}, 2026.

  \item Kit Pagos Colombia, ``Documentación de datos de prueba,'' \path{docs/testing-data/}, 2026.

  \item Kit Pagos Colombia, ``Prototipos del experimento de la Fase 5,'' \path{prototypes/README.md}, 2026.

  \item Kit Pagos Colombia, ``Colección de Postman de la API de Simulación,'' \path{simulator-api/postman/}, 2026.

  \item Meta Open Source, ``Jest: JavaScript testing framework.'' [En línea]. Disponible en: \url{https://jestjs.io/docs/getting-started}

  \item Fastify, ``Fastify documentation.'' [En línea]. Disponible en: \url{https://fastify.dev/docs/latest/}

  \item Postman, ``Postman learning center.'' [En línea]. Disponible en: \url{https://learning.postman.com/docs/}

  \item OpenAPI Initiative, ``OpenAPI Specification v3.0.3,'' 2021. [En línea]. Disponible en: \url{https://spec.openapis.org/oas/v3.0.3}

  \item GitHub, ``GitHub Actions documentation.'' [En línea]. Disponible en: \url{https://docs.github.com/en/actions}

  \item OpenJS Foundation, ``Node.js releases.'' [En línea]. Disponible en: \url{https://nodejs.org/en/about/previous-releases}

  \item Wompi, ``Documentación para desarrolladores.'' [En línea]. Disponible en: \url{https://docs.wompi.co/}

  \item Mercado Pago, ``Documentación para desarrolladores.'' [En línea]. Disponible en: \url{https://www.mercadopago.com.co/developers/es/docs}

  \item Kushki, ``Documentación para desarrolladores.'' [En línea]. Disponible en: \url{https://docs.kushkipagos.com/}

  \item Rapyd, ``Rapyd API documentation.'' [En línea]. Disponible en: \url{https://docs.rapyd.net/}

  \item M. Collina \emph{et al.}, ``autocannon: fast HTTP/1.1 benchmarking tool written in Node.js,'' versión 8.0.0. El histograma de latencia se define en el archivo \emph{histUtil.js} de la carpeta \emph{lib}, con un valor mínimo distinguible de 1~ms, y \emph{run.js}, en la misma carpeta, solo registra las latencias de respuestas 2xx. [En línea]. Disponible en: \url{https://github.com/mcollina/autocannon}\label{ref:autocannon}

  \item Fastify, ``Fastify benchmarks.'' [En línea]. Disponible en: \url{https://github.com/fastify/benchmarks}\label{ref:fastify-benchmarks}

  \item OpenJS Foundation, ``Performance measurement APIs (perf\_hooks),'' documentación de Node.js. [En línea]. Disponible en: \url{https://nodejs.org/api/perf_hooks.html}\label{ref:node-perf-hooks}

  \item D. Eastlake y A. Panitz, ``Reserved Top Level DNS Names,'' RFC 2606, 1999. [En línea]. Disponible en: \url{https://www.rfc-editor.org/rfc/rfc2606.txt}\label{ref:rfc2606}

  \item IEEE, ``IEEE Standard for Software and System Test Documentation,'' IEEE Std 829-2008, 2008.

  \item PCI Security Standards Council, ``PCI DSS: Payment Card Industry Data Security Standard.'' [En línea]. Disponible en: \url{https://www.pcisecuritystandards.org/}
\end{enumerate}

\section{Anexos}
\label{sec:anexos}

Esta sección reúne el material complementario que respalda la ejecución del plan de pruebas: la matriz de trazabilidad completa entre requisitos y casos, la guía de comandos para la corrida de verificación y el glosario de términos técnicos.

\subsection{Matriz de trazabilidad}

La Tabla~\ref{tab:trazabilidad} relaciona los requisitos funcionales e historias de usuario del SRS con los casos de este plan. Un requisito sin caso es una brecha declarada; un caso sin requisito es una verificación que no nace de una necesidad del producto.

{\small
\begin{longtable}{|>{\raggedright\arraybackslash}p{1.3cm}|>{\raggedright\arraybackslash}p{7.7cm}|>{\raggedright\arraybackslash}p{2.6cm}|>{\raggedright\arraybackslash}p{3.2cm}|}
\caption{Matriz de trazabilidad entre requisitos y casos}\label{tab:trazabilidad}\\
\hline
\textbf{ID} & \textbf{Requisito} & \textbf{Casos} & \textbf{Historias} \\ \hline
\endfirsthead
\caption[]{Matriz de trazabilidad entre requisitos y casos (continuación)}\\
\hline
\textbf{ID} & \textbf{Requisito} & \textbf{Casos} & \textbf{Historias} \\ \hline
\endhead
\hline
\endfoot

RF-01 & Crear una intención de pago con tarjeta o PSE mediante una interfaz unificada. & CP-01, CP-02, CP-04, CP-05 & HU-01 \\ \hline
RF-02 & Normalizar la respuesta a una \texttt{Transaction} con estado y motivo. & CP-01, CP-03, CP-07 & HU-01, HU-02 \\ \hline
RF-03 & Consultar el estado y mapearlo a los seis estados normalizados. & CP-01, CP-02, CP-07 & HU-02 \\ \hline
RF-04 & Validar la firma criptográfica de un webhook. & CP-04 & HU-03 \\ \hline
RF-05 & Convertir los fallos esperados en \texttt{KitPagosError}. & CP-03, CP-07 & HU-05, HU-11 \\ \hline
RF-06 & Cambiar de pasarela sin cambiar el código de negocio. & CP-01, CP-04 & HU-04 \\ \hline
RF-07 & Reintentar operaciones idempotentes con retroceso exponencial. & CP-07 & HU-05 \\ \hline
RF-08 & Redactar credenciales y secretos de los errores. & CP-03, CP-07 & HU-05 \\ \hline
RF-09 & Exponer las rutas nativas bajo \texttt{/v1/sim}. & CP-01 a CP-05, CP-07 & HU-06, HU-07 \\ \hline
RF-10 & Elegir el escenario por cabecera, dato, monto o marca. & CP-02, CP-03, CP-07 & HU-06, HU-07 \\ \hline
RF-11 & Replicar la estructura de error nativa. & CP-03, CP-07 & HU-07 \\ \hline
RF-12 & Disparar y emitir webhooks firmados. & CP-04 & HU-08 \\ \hline
RF-13 & Publicar la especificación OpenAPI 3.0. & CP-06 & HU-09 \\ \hline
RF-14 & Documentar datos de prueba por pasarela. & CP-02, CP-03, CP-05 & HU-13 \\ \hline
RF-15 & Indicar el comportamiento y la fuente de cada dato. & CP-03, CP-05 & HU-13 \\ \hline
RF-16 & Exponer las operaciones por REST bajo \texttt{/v1/api}. & CP-06 & HU-09 \\ \hline
RF-17 & Tokenizar la tarjeta en el navegador. & CP-05 & HU-10 \\ \hline
RF-18 & Exponer el motivo de rechazo como código y categoría. & CP-03, CP-06 & HU-11 \\ \hline
RF-19 & Listar los bancos PSE y traducir su código. & CP-01, CP-02, CP-06 & HU-12 \\ \hline
\end{longtable}
}

\subsection{Comandos de la corrida}

La Tabla~\ref{tab:comandos} reúne los comandos que ejecutan las suites y los ejemplos. Se escriben desde la raíz del repositorio salvo que la fila indique otra carpeta.

\begin{table}[H]
  \centering
  \small
  \renewcommand{\arraystretch}{1.2}
  \setlength{\tabcolsep}{3pt}
  \begin{tabularx}{\textwidth}{|>{\raggedright\arraybackslash}p{3.2cm}|>{\raggedright\arraybackslash}X|>{\raggedright\arraybackslash}p{4.2cm}|}
    \hline
    \textbf{Comprobación} & \textbf{Comando} & \textbf{Resultado esperado} \\ \hline
    Unitarias del SDK & \texttt{cd sdk \&\& npm test -- --coverage} & 58 suites y 1088 pruebas en verde. \\ \hline
    Contrato & \texttt{cd sdk \&\& npm run test:sandbox} & 25 pruebas ejecutadas o omitidas por falta de \texttt{.env}. \\ \hline
    API de Simulación & \texttt{cd simulator-api \&\& npm test -- --coverage} & 54 suites y 1301 pruebas en verde. \\ \hline
    Escenarios & \texttt{cd examples \&\& npm run simulate:scenarios} & Cada línea con \texttt{OK}; código de salida 0. \\ \hline
    Intercambiabilidad & \texttt{cd examples \&\& npm run simulate:interchangeability} & Cuatro pasarelas con el mismo resultado; código 0. \\ \hline
    PSE de Wompi & \texttt{cd examples \&\& npm run simulate:wompi-pse} & Redirección \texttt{PENDING} y consulta \texttt{APPROVED}. \\ \hline
    Typecheck de ejemplos & \texttt{cd examples \&\& npm run typecheck} & Compila contra la superficie pública del SDK. \\ \hline
    Prototipos & \texttt{npm test -- --coverage} y \texttt{npm run typecheck} en cada carpeta & Suites y compilación en verde. \\ \hline
    Paridad de esqueleto & \texttt{./check-skeleton-parity.sh} & Sin diferencias entre B y D. \\ \hline
    Postman local & Importar la colección con el ambiente \texttt{local} & 42 peticiones con sus aserciones en verde. \\ \hline
    Postman Render & Importar la colección con el ambiente \texttt{render} & 42 peticiones con sus aserciones en verde. \\ \hline
    OpenAPI & \texttt{GET /docs}, \texttt{/docs/json} y \texttt{/docs/yaml} & Las tres rutas responden y el documento describe las rutas reales. \\ \hline
  \end{tabularx}
  \caption{Comandos de la corrida}
  \label{tab:comandos}
\end{table}

\subsection{Glosario mínimo}

\begin{description}
  \item[Sandbox.] Ambiente de pruebas de una pasarela real. Acepta credenciales de prueba y puede declinar sin motivo aparente.
  \item[Simulador.] API de Simulación local o desplegada en Render. Es determinista, no necesita credenciales reales y replica la forma de las rutas nativas.
  \item[Token opaco.] Identificador que reemplaza los datos de la tarjeta. El comercio lo recibe del navegador o de la pasarela y no puede interpretarlo.
  \item[Estado normalizado.] Uno de los seis valores de \texttt{TransactionStatus}: \texttt{APPROVED}, \texttt{DECLINED}, \texttt{PENDING}, \texttt{EXPIRED}, \texttt{VOIDED} o \texttt{ERROR}.
  \item[Estado nativo.] El valor que devuelve la pasarela (\texttt{rawStatus}), conservado para auditoría.
  \item[Paridad de esqueleto.] Comprobación de que los prototipos B y D comparten servidor, contrato, webhooks, configuración y dependencias; solo cambia el módulo de pagos.
\end{description}


\end{document}
